% Préambule repris du fichier officiel GabaritAMQ101.tex téléchargé le 23 août 2026.
% Les dimensions de page et les marges officielles sont conservées sans modification.
\usepackage[T1]{fontenc}
\usepackage{lmodern}
\usepackage[utf8]{inputenc}
\usepackage[english,french]{babel}
\usepackage{graphicx}
\usepackage{amssymb,latexsym,amsmath}
\usepackage{mathrsfs}
\usepackage{multirow}
\usepackage{euscript}
\usepackage{fancyhdr}
\usepackage{pdfsync}
\usepackage{theorem}
\usepackage{wrapfig}
\usepackage{floatflt}
\usepackage[protrusion=true]{microtype}
\usepackage[all]{xy}
\usepackage{geometry}
\usepackage{tikz}
\usepackage{booktabs}
\usepackage{color}
\usepackage{url}
\usepackage[colorlinks=true,urlcolor=blue,citecolor=cyan,linkcolor=magenta]{hyperref}
\usepackage{esvect}

\renewcommand{\vec}[1]{\protect\vv{#1}}

\newtheorem{theorem}{Théorème}
\newtheorem{lemma}{Lemme}
\newtheorem{proposition}{Proposition}
\newtheorem{corollary}{Corollaire}
\newtheorem{definition}{Définition}
\newtheorem{nota}{Notation}
\newtheorem{ex}{Exemple}
\newtheorem{nb}{N.B.}
\newtheorem{remark}{Remarque}

\newenvironment{proof}
{\par\noindent\textit{Démonstration}\ }
{\hfill{$\Box$}}

\newcommand{\siecle}[1]{%
\ifnum #1=1%
\uppercase\expandafter{\romannumeral #1}\textsuperscript{er}%
\else%
\uppercase\expandafter{\romannumeral #1}\up{e}%
\fi}

% Définition des marges pour le Bulletin, à ne pas modifier.
\geometry{height=18.6cm,width=14.4cm}
\geometry{hmargin={3.6cm,3.6cm}}
\geometry{vmargin={4.3cm,5cm}}
\parindent=0.0in
\parskip=0.1in
\setcounter{page}{1}

\DeclareSymbolFont{AMSb}{U}{msb}{m}{n}
\DeclareSymbolFontAlphabet{\Bbb}{AMSb}
\newcommand{\ds}{\displaystyle}

\usepackage[]{lineno}
\newcommand*\patchAmsMathEnvironmentForLineno[1]{%
  \expandafter\let\csname old#1\expandafter\endcsname\csname #1\endcsname
  \expandafter\let\csname oldend#1\expandafter\endcsname\csname end#1\endcsname
  \renewenvironment{#1}%
     {\linenomath\csname old#1\endcsname}%
     {\csname oldend#1\endcsname\endlinenomath}}
\newcommand*\patchBothAmsMathEnvironmentsForLineno[1]{%
  \patchAmsMathEnvironmentForLineno{#1}%
  \patchAmsMathEnvironmentForLineno{#1*}}
\AtBeginDocument{%
\patchBothAmsMathEnvironmentsForLineno{equation}%
\patchBothAmsMathEnvironmentsForLineno{align}%
\patchBothAmsMathEnvironmentsForLineno{flalign}%
\patchBothAmsMathEnvironmentsForLineno{alignat}%
\patchBothAmsMathEnvironmentsForLineno{gather}%
\patchBothAmsMathEnvironmentsForLineno{multline}%
}

\usepackage{url}
\usepackage{hyperref}
\newcommand*\rfrac[2]{{}^{#1}\!/_{#2}}
\newcommand{\fracinline}[2]{\raisebox{0.4ex}{$#1$} / \raisebox{-0.7ex}{$#2$}}
\newcommand{\bigfracinline}[2]{\raisebox{0.8ex}{$#1$} \Big/ \raisebox{-1.4ex}{$#2$}}
\newcommand{\biggfracinline}[2]{\raisebox{0.8ex}{$#1$} \Bigg/ \raisebox{-1.4ex}{$#2$}}
\everymath{\displaystyle}
